\section{Introduction}
\label{sec:intro}

Robot manipulation policies are converging on two ways of turning pixels and
language into motor commands. The first family starts from a vision-language
model and attaches an action head to it, so the policy inherits semantic breadth
from web-scale pretraining and learns the map from an instruction to a motion
directly. The second family starts from a video generative model and asks a
single network to produce the future the robot is about to see together with the
actions that produce it, which places a model of the world inside the policy
itself. This report studies the second family, which the literature calls
world-action models. Two properties make that family compelling. A video
objective supplies dense supervision from footage whose action labels are
absent, so the policy acquires physical regularity from human activity and from
unannotated robot logs. The predicted future then acts as an intermediate
quantity that ties every candidate action to a visible consequence, which
grounds control in the same representation that already carries scene dynamics.

Building a world-action model requires a long sequence of independent decisions,
and every published system settles all of them at once. A designer selects the
video generator that supplies the backbone. The designer selects the visual
representation that this generator denoises. The designer selects the way the
action stream attaches to the video stream, and the amount of the video
computation that the action stream is permitted to read. The designer selects
the corpus the joint model is pretrained on and the order in which the two
streams are denoised once the policy drives a robot. Each published result
reports one point of this space together with one benchmark score. These
decisions interact strongly, so the reasoning that justifies a choice inside one
system loses its force inside another. A designer who sets out to build a new
world-action model therefore inherits a collection of working systems together
with an open question about which choice actually carries the performance.

We address this by treating the design space itself as the object of study. The
report contributes a framework in which each decision above becomes an
independent component behind a fixed interface. One configuration file selects
the video backbone and the visual representation it denoises. The same file
selects how the action stream couples to that backbone and how much of the
backbone the action stream reads. The same file selects the pretraining mixture
and the schedule that drives the policy at run time. The framework is the
instrument that makes measurement possible, because interchangeable components
allow a study to hold an entire configuration fixed and move one axis at a time
under one dataset, one optimisation budget, one evaluation harness and one
metric. We name the framework and the model family it produces \sys, for joint
action-world modelling. The pretraining it performs is action grounded in a
precise sense, because one denoising process predicts the future observation and
predicts the actions that bring that observation about, and the visibility
pattern between the two streams sets how far each one informs the other.

The study asks six questions in the order a designer meets them. It asks how the
action stream should be coupled to the video stream. It asks which video
generator repays its computational cost. It asks what property of a visual
representation makes that representation a good denoising target for a policy.
It asks how much of the video computation the action stream should read, and
whether large-scale pretraining moves the answer. It asks which mixture of human
egocentric footage and robot trajectories to pretrain on, and whether the two
sources should be staged or combined. It asks how the two streams should be
denoised once the policy controls a robot. We then compose the configuration
that the study selects, pretrain it on a large mixture of human egocentric
footage and robot trajectories spanning simulation and real hardware, and
fine-tune the result for every benchmark and every physical platform we
evaluate.

The value of this report follows from the completeness of what it releases. Every
point of the study corresponds to a checkpoint we publish, so a reader who
doubts a conclusion can reload the exact model that produced it. The pretrained
model demonstrates that the resulting map is actionable, because that model is
precisely the configuration the study selected, scaled to a large corpus and
trained once. A designer consequently gains a defensible default for every
decision, an account of which decisions deserve additional capacity, and an open
path that runs from raw footage to a policy executing on physical hardware. Our
contributions follow this structure. We build a modular framework that turns the
world-action design space into a set of interchangeable components sharing one
training, deployment and evaluation interface (\sref{sec:framework}). We run a
controlled six-axis study of that space under a single protocol
(\sref{sec:study}). We release a pretrained foundation model together with
fine-tuned variants covering a broad suite of simulated benchmarks and several
real robots (\sref{sec:alpha}).
